\documentclass[a4paper,11pt]{ltjsarticle}
\usepackage[haranoaji]{luatexja-preset}
\usepackage{amsmath,amssymb}
\usepackage[top=12mm,bottom=10mm,left=14mm,right=14mm]{geometry}
\usepackage{array}
\usepackage{tikz}
\usetikzlibrary{calc}
\pagestyle{empty}
\setlength{\parindent}{0pt}
\newlength{\qcol}\setlength{\qcol}{0.47\textwidth}
\newlength{\qgap}
\newlength{\anscolw}\setlength{\anscolw}{0.30\textwidth}
\newlength{\ansh}
\newcommand{\Q}[2]{\parbox[t]{\qcol}{\raggedright (#1)\hspace{0.6em}$\displaystyle #2$}}
\newcommand{\QR}[4]{\Q{#1}{#2}\hfill\Q{#3}{#4}\par\vspace{\qgap}}
\newcommand{\anscell}[1]{\rule[-\ansdrop]{0pt}{\ansh}(#1)}
\newlength{\ansdrop}

\begin{document}
{\large\bfseries 計算問題　21}\hfill 名前(\underline{\hspace{32mm}})\par\vspace{3mm}
■（1）〜（16）の計算をしなさい。（17）〜（21）は方程式を解きなさい。\par\vspace{3mm}
\setlength{\qgap}{5.4mm}
\QR{1}{(-21)\div(+3)}{2}{(-14)+(-6)-(-7)}
\QR{3}{8\times\{3-(-2+1)\times2\}}{4}{\dfrac{7}{3}(3x-15)-\dfrac{2}{5}(5x+35)}
\QR{5}{(21a-6)\div3}{6}{\dfrac{7}{3}\div2+\dfrac{5}{3}}
\QR{7}{(-11)\times(-5)}{8}{-7^2+(-3)^2-2}
\QR{9}{15-(-3)\div\left(-\dfrac{3}{2}\right)}{10}{\dfrac{1}{7}(21x-14)+\dfrac{1}{5}(10x+15)}
\QR{11}{(3x-7)-(5x+2)}{12}{\dfrac{x-7}{3}\times6}
\QR{13}{(7a+3)+(2a-5)}{14}{7\div3\times(-6)}
\QR{15}{7(3x-2)-(2x-5)}{16}{-0.7-(-2.3)-1.5}
\QR{17}{7x-3=-38}{18}{7(x+3)=8x+24}
\QR{19}{0.7x+0.3=-2.5}{20}{\dfrac{3}{4}x-2=\dfrac{1}{2}x+1}
\vspace{1mm}
\Q{21}{7(x-2)-3(x+1)=10x-13}\par\vspace{3mm}
\setlength{\ansh}{9.5mm}\setlength{\ansdrop}{6mm}%
\begin{center}\begin{tabular}{|p{\anscolw}|p{\anscolw}|p{\anscolw}|}\hline
\anscell{1} & \anscell{2} & \anscell{3} \\\hline
\anscell{4} & \anscell{5} & \anscell{6} \\\hline
\anscell{7} & \anscell{8} & \anscell{9} \\\hline
\anscell{10} & \anscell{11} & \anscell{12} \\\hline
\anscell{13} & \anscell{14} & \anscell{15} \\\hline
\anscell{16} & \anscell{17} & \anscell{18} \\\hline
\anscell{19} & \anscell{20} & \anscell{21} \\\hline
\end{tabular}\end{center}

\end{document}
